\section{Evaluation}

% claims: C08; sources: S10 S11
We organise evaluation around five scientific questions. The primary unit for
the integrated tests is the closure defined in Section~2: one matched
natural/alias pair judged by an all-or-nothing endpoint. Generated outputs,
positions, authority scores, receipt checks, and attack checks inside a closure
are dependent measurements.

\paragraph{Q1: Is the internal state causally operative?}
% claims: C02 C03; sources: S03 S04
We intervene on hard \ztwo{} while leaving the learned transition responsible
for \zthree{}. Alias interventions must change \ztwo{} and \zthree{} while
preserving the answer; non-alias controls must change the answer. Same-state,
sham, wrong-state, and direct-\zthree{} families test target consumption and
registered downstream behaviour. This is a prospective selected-model test,
not evidence that the state emerged naturally or that the recipe is robust
across seeds.

\paragraph{Q2: Can exact and abstract evidence be authenticated?}
% claims: C04 C05; sources: S05 S06
Genuine exact, abstract, and chain receipts are evaluated together with frozen
mutation, replay, crossing, omission, and context-substitution tests. The
endpoint asks whether the trusted verifier distinguishes a precise event, a
reusable state class, and the registered causal chain. It does not test a
malicious issuer, compromised key, or public nonrepudiation.

\paragraph{Q3: Can the carrier transmit a generic authority?}
% claims: C06 C07; sources: S07 S08 S09
Clean carrier-off calibration fixes the threshold and margin before a separate
prospective test of 16 generic authorities. The generic test uses 512 closures
and no receipts, isolating statistical carriage from state authentication. It
tests the frozen operating point in this grammar, not channel capacity,
imperceptibility, or robustness to editing.

\paragraph{Q4: Does the integrated signal follow authenticated state?}
% claims: C08 C09; sources: S10 S11 S12 S13
The public integrated evaluation uses 128 closures. In each, natural and alias
executions share the prompt, answer, surface payload, configuration, output
indices, and indexed random draws; authenticated \ztwo{} and its registered
downstream consequences switch as in Figure~\ref{fig:contrast}. A closure passes
only if both sides satisfy computation fidelity; exact, abstract, and chain
verification; parsing; intended-authority acceptance; carrier-off rejection;
absence of wrong-authority confusion; and every frozen protocol-consistency
attack. The comparison controls these listed alternatives within the trusted
pipeline; it cannot rule out an arbitrary controller already trusted with the
state label.

\paragraph{Q5: Does the result survive public and protected evaluation?}
% claims: C08 C10 C11; sources: S10 S11 S14 S15 S16 S17
Public validation, protected qualification, and protected confirmation have
separate populations and frozen roles. Gates are non-compensatory: a carrier
success cannot offset a computation or attack failure. ``Held out'' means that
public closures were disjoint from model training and released to integration
only after calibration and generic validation were frozen; it does not mean
independent models, grammars, or templates. Public and protected populations
share the checkpoint, ontology, grammar, and controller but use distinct
closure identities and keys. Their outcomes must not be pooled.
